\documentclass[11pt,a4paper]{article}
\newcommand{\labonaamwaarde}{lab 5}
\usepackage[english]{syllabus}

\begin{document}

\syllabustitel{lab 5 -- \emph{Data Visualization I}}{Tell the story with plotly}

% ---------------------------------------------------------------------------
% Shared by the syllabi of labs 3, 4 and 5. Edit it here, not in the lab files.
\section{Overview}

In labs 3, 4 and 5 you are the new data analyst of a (fictional) youth
football and athletics club. The head coach gives you three export files
from the club's systems. They are messy on purpose, the way real exports
are: inconsistent labels, sprint times in milliseconds, missing values,
an athlete who appears twice and a column with a comma as the decimal
separator.

Each session builds on the previous one:

\begin{description}[leftmargin=1.6cm, style=nextline]
  \item[Lab 3] \textbf{From file to tidy table.} Load, explore, filter,
    add columns and clean until you can trust the data.
  \item[Lab 4] \textbf{Combine, summarize, reshape.} Summarize groups,
    combine tables, switch between a wide and a long shape, and work with
    a heart-rate series per second.
  \item[Lab 5] \textbf{Tell the story with plotly.} Choose the right chart
    for a question, make honest charts with \code{plotly.express} and
    combine them into a small interactive dashboard.
\end{description}


% ---------------------------------------------------------------------------
\section{Learning outcomes}

After lab 5 you can:

\begin{enumerate}[label=\textbf{LO\arabic*}, leftmargin=1.4cm]
  \item choose a chart type that fits the question, and recognise
        misleading charts, such as a truncated axis;
  \item make histograms, box plots, scatter plots and line charts with
        \code{plotly.express}, with colour, facets and size for extra
        dimensions;
  \item make charts interactive with marimo UI elements and combine them
        into a dashboard.
\end{enumerate}

% ---------------------------------------------------------------------------
\section{Prerequisites}

Labs 3 and 4 (\emph{Data Processing I and II}): loading, filtering and
cleaning data with \code{pandas}, grouping with \code{groupby}, combining
tables with \code{merge}, and the difference between a long and a wide
shape. Lab 5 starts with the complete cleaning of labs 3 and 4 in one
cell.

% ---------------------------------------------------------------------------
% Shared by the syllabi of labs 3, 4 and 5. Edit it here, not in the lab files.
\section{Software and materials}

Every lab is a \textbf{marimo notebook} in the \code{labs/} folder of the
course repository. You can open it in two ways:

\begin{itemize}
    \item \textbf{Online:} go to \url{https://marimo.app}, choose
        \emph{New $\rightarrow$ Open from URL\dots} and paste the raw URL
        of the notebook on GitHub. The data is then read from GitHub, so
        you don't have to install anything.
    \item \textbf{Locally:} install the packages from
        \code{requirements.txt} and open the notebook with marimo. The
        data is then read from the \code{data/} folder.
\end{itemize}

Three things to know about marimo:

\begin{itemize}
  \item Notebooks are \textbf{reactive}: when you change a cell, every
        cell that depends on it runs again, in the order of their
        dependencies.
  \item \textbf{Two cells cannot define the same variable.} Give every
        step its own clear name, such as \code{tests\_raw} and
        \code{tests\_clean}, or keep a sequence of steps together in one
        cell.
  \item Every exercise contains a \code{...} placeholder. A hidden cell
        below it checks your answer: it turns \textbf{green} when the
        answer is right, and otherwise gives a general hint about what is
        still wrong, without giving the answer away.
\end{itemize}


% ---------------------------------------------------------------------------
% Shared by the syllabi of labs 3, 4 and 5. Edit it here, not in the lab files.
\section{The dataset}

\noindent
\rowcolors{2}{tint}{white}
\begin{tabularx}{\textwidth}{@{}L{3.1cm}L{3.2cm}Y@{}}
\toprule
\textbf{File} & \textbf{Size} & \textbf{Content and known problems} \\
\midrule
\code{athletes.csv}   & 43 rows, 7 columns   & ID, name, birthdate, gender, position, height, weight. 13 spellings for 4 positions, weight as text with a decimal comma (\code{"57,5"}), one athlete entered twice with two IDs, a few missing values. \\
\code{tests.csv}      & 114 rows, 8 columns  & 10\,m and 30\,m sprint, countermovement jump height and Yo-Yo level, at 3 test moments. Some sprint times are stored in milliseconds; a few missing values. \\
\code{hr\_session.csv} & 12,000 rows, 5 columns & Heart rate and speed per second for 5 athletes during one 40-minute training session. \\
\bottomrule
\end{tabularx}
\rowcolors{1}{}{}


% ---------------------------------------------------------------------------
\needspace{18\baselineskip}
\section{Schedule}

\textbf{Goal:} choose a chart type for a given question, make the
standard \code{plotly.express} charts, add extra dimensions with colour,
facets and size, and build in simple interactivity.

\begin{schedule}
13:00--13:30 & Principles & The right chart for the question (distribution, comparison, relationship, evolution); critique a misleading chart. \\
13:30--14:30 & \code{plotly.express} basics & Histogram, box plot, scatter plot and line chart, with titles, axis labels and hover information. \\
14:30--15:15 & More dimensions & \code{color}, \code{facet\_col} and \code{size}; one fixed colour per position. \\
15:15--15:30 & Break & \\
15:30--16:15 & Heart rate as a time series & Zone bands with \code{add\_hrect}, and why we don't use a second y-axis. \\
16:15--16:45 & Interactivity & \code{mo.ui.dropdown}, \code{mo.ui.range\_slider}, optionally selections with \code{mo.ui.plotly}. \\
16:45--17:30 & Mini-dashboard & 3 charts and 1 selector that at least one chart responds to. \\
\end{schedule}

\textbf{Exercises:} (1) a histogram of weight; (2) Yo-Yo level per
position, with facets per gender; (3) speed over time for one athlete;
\textbf{assignment:} the mini-dashboard.

% ---------------------------------------------------------------------------
\section{Working method}

\begin{itemize}
  \item Each block starts with a short demo in the notebook, followed by
        exercises you solve in the notebook itself.
  \item Use the automatic checks as feedback, not as the goal: when a
        check turns green, make sure you can explain \emph{why} your code
        works.
  \item You build the mini-dashboard with everything from this lab: the
        goal is combining, not learning new chart types.
\end{itemize}

% ---------------------------------------------------------------------------
\section{Assessment}

\begin{itemize}
  \item \textbf{Mini-dashboard:} three charts and one interactive
        selector, made with the club's data. It can serve as your
        submission for this module and prepares the case presentations
        and the Project Lab next semester.
\end{itemize}

% ---------------------------------------------------------------------------
\section{Preparation}

\begin{itemize}
  \item Check before the session that you can open the notebook, locally
        or on marimo.app.
  \item Review \code{groupby}, \code{merge} and the long shape from lab 4.
        Charts with plotly express start from data in the long shape.
  \item Read the course text for lab 5 next to the notebook.
\end{itemize}

\end{document}
